% ===== PRÉAMBULE CANONIQUE — à coller VERBATIM en tête de CHAQUE .tex =====
\documentclass[11pt,a4paper]{article}
\usepackage[utf8]{inputenc}\usepackage[T1]{fontenc}\usepackage{lmodern}
\usepackage[french]{babel}
\usepackage{amsmath,amssymb,mathtools}\usepackage{icomma}
\usepackage{siunitx}\sisetup{locale=FR}  % \qty{}{}, \si{}, \ang{}
\usepackage{xcolor}\usepackage{colortbl}
\usepackage{tikz}\usepackage{pgfplots}\pgfplotsset{compat=1.18}
\usepackage[siunitx,europeanresistors]{circuitikz} % circuits (élec.)
\usepackage[most]{tcolorbox}\usepackage{enumitem}\usepackage{array,booktabs}
\usepackage[a4paper,margin=2.2cm]{geometry}
\usetikzlibrary{arrows.meta,calc,angles,quotes,positioning,patterns,%
  backgrounds,shapes.geometric,decorations.pathmorphing,decorations.markings,intersections}
\definecolor{primary}{RGB}{222,49,99}\definecolor{primarydark}{RGB}{150,22,55}
\definecolor{defcol}{RGB}{0,114,178}\definecolor{methcol}{RGB}{52,92,148}
\definecolor{excol}{RGB}{0,135,90}\definecolor{attncol}{RGB}{200,40,55}
\tcbset{boxbase/.style={enhanced,breakable,boxrule=0.9pt,arc=2.5pt,
  left=3mm,right=3mm,top=2.3mm,bottom=2.3mm,fonttitle=\bfseries,coltitle=white}}
% --- boîtes TUTORIEL (noms EXACTS, ne pas renommer) ---
\newtcolorbox{cle}[1][]{boxbase,colback=defcol!6,colframe=defcol,
  title={\textsc{À retenir}\ifx\relax#1\relax\relax\else\textmd{ : \itshape #1}\fi}}
\newtcolorbox{methode}[1][]{boxbase,colback=methcol!7,colframe=methcol,sharp corners,
  title={\textsc{Méthode}\ifx\relax#1\relax\relax\else\textmd{ --- \itshape #1}\fi}}
\newtcolorbox{exemplebox}[1][]{boxbase,colback=excol!5,colframe=excol,
  title={\textsc{Exemple}\ifx\relax#1\relax\relax\else\textmd{ : \itshape #1}\fi}}
\newtcolorbox{piege}[1][]{boxbase,colback=attncol!6,colframe=attncol,
  title={\textsc{Piège}\ifx\relax#1\relax\relax\else\textmd{ : \itshape #1}\fi}}
\newtcolorbox{remarque}[1][]{boxbase,colback=black!4,colframe=black!45,
  title={\textsc{Remarque}\ifx\relax#1\relax\relax\else\textmd{ : \itshape #1}\fi}}
% --- boîtes EXERCICE / CORRIGÉ (noms EXACTS, auto-comptées) ---
\newtcolorbox[auto counter]{exo}[1][]{boxbase,colback=primary!4,colframe=primary,
  title={\textsc{Exercice}~\thetcbcounter\ifx\relax#1\relax\relax\else\textmd{ --- \itshape #1}\fi}}
\newtcolorbox[auto counter]{corrige}[1][]{boxbase,colback=excol!4,colframe=excol,
  title={\textsc{Corrigé}~\thetcbcounter\ifx\relax#1\relax\relax\else\textmd{ --- \itshape #1}\fi}}
\newcommand{\vv}[1]{\vec{#1}}\newcommand{\ds}{\displaystyle}
\newcommand{\R}{\mathbb{R}}\newcommand{\N}{\mathbb{N}}\newcommand{\Z}{\mathbb{Z}}
% =========================================================================
\begin{document}
\sisetup{per-mode=symbol}

\begin{corrige}[deuxième loi : $a=F/m$]
On applique le PFD, $F=ma$, projeté sur la direction du mouvement.
\begin{enumerate}[label=\alph*)]
  \item $a=\dfrac{F}{m}=\dfrac{12}{3}=\qty{4}{\metre\per\second\squared}$.
  \item L'accélération est \emph{proportionnelle} à la force : en doublant $F$, on double $a$, soit
        $a=\dfrac{24}{3}=\qty{8}{\metre\per\second\squared}$.
\end{enumerate}
\end{corrige}

\begin{corrige}[deuxième loi : $F=ma$]
\begin{enumerate}[label=\alph*)]
  \item $F=m\,a=2\times 5=\qty{10}{\newton}$.
  \item $[F]=\si{\kilogram}\times\si{\metre\per\second\squared}=\si{\kilogram\metre\per\second\squared}=\si{\newton}$ : cohérent.
\end{enumerate}
\end{corrige}

\begin{corrige}[retrouver la masse]
De $F=ma$ on tire $m=\dfrac{F}{a}=\dfrac{50}{\num{2,5}}=\qty{20}{\kilogram}$.
\end{corrige}

\begin{corrige}[première loi : la résultante]
\begin{enumerate}[label=\alph*)]
  \item Les deux forces se compensent exactement : $\vv N+\vv G=\vv 0$, donc $\vv F_{\text{rés}}=\vv 0$.
  \item D'après la $1^{\text{re}}$ loi, $\vv F_{\text{rés}}=\vv 0 \Rightarrow \vv a=\vv 0$ : le livre reste
        \emph{au repos} (il pourrait aussi être en MRU, mais ici sa vitesse initiale est nulle).
\end{enumerate}
\end{corrige}

\begin{corrige}[direction de l'accélération]
\begin{enumerate}[label=\alph*)]
  \item $\vv a$ est \emph{colinéaire} à la résultante : même direction (horizontale) et même sens
        (vers la droite) que $\vv F$.
  \item Non : $\vv F$ étant l'unique force, $\vv a=\vv F/m$ a forcément la direction et le sens de $\vv F$.
        Une accélération vers le haut exigerait une force résultante vers le haut.
\end{enumerate}
\end{corrige}
\end{document}
